\ncertchapter{8}{Connected Bodies: Pulleys, Strings and Springs}\chlabel{ch:connected}

A single body under known forces is a straightforward application of $\vF=m\va$. Most interesting
mechanical systems, however, consist of several bodies joined together: two masses over a pulley, a block
on a table dragged by a hanging weight, a load supported by a spring through a system of strings. Joining
bodies introduces unknowns, the tensions in the strings and the forces in the springs, and the question is
how to find them along with the motion.

The answer rests on two ideas, and the chapter is organised around them. The first is that an ideal string,
one that is light and inextensible, carries a single tension throughout its length and imposes a
\emph{constraint} linking the accelerations of the bodies it connects. Writing that constraint correctly,
by asking how the string's length is shared, is the step students most often get wrong. The second idea is that the networks of light springs met in practice
reduce to a single effective spring, and that the rule for doing so depends on what the springs share: the
same force, or the same extension.

Throughout, strings are ideal (light and inextensible) and pulleys are light and frictionless unless a
problem says otherwise; \chref{ch:newton} showed that these are exactly the conditions under which one
string carries one tension.

\section{The Standard Machines}

The three classic connected-body problems (the Atwood machine, the block on a table and the block on an
incline) are really \emph{one} machine seen at different tilts, so one method solves all three and saves
a great deal of memorisation. In each, the string supplies both a force (the tension) and a constraint
linking the accelerations. If each body's own direction of motion is taken as positive, the tension,
which pulls each body towards the pulley, enters the two equations with opposite signs and cancels when
they are added.

\subsection{The Atwood machine}

The simplest connected system is a pair of unequal masses hanging over a fixed pulley, and the
method that solves it also solves every other system in the chapter. George Atwood built it in 1784 to slow down free
fall: with nearly equal masses the acceleration is a small fraction of $g$, slow enough to time by hand.

\begin{figure}[htbp]
\centering
\tikzfig{atwood}
\figcaption{fig:atwood}{The Atwood machine with $m_1>m_2$. The same tension $T$ acts on both masses, and it lies between the two weights: smaller than $m_1g$, so $m_1$ accelerates down, and larger than $m_2g$, so $m_2$ accelerates up, both with the same magnitude $a$.}
\end{figure}

\begin{derivbox}{The Atwood machine}
Two masses $m_1>m_2$ hang from the ends of a light, inextensible string over a fixed pulley that is
light and frictionless. The heavier one falls and the
lighter rises; because the string is inextensible, both move with the \emph{same} acceleration magnitude
$a$, and the tension $T$ is the same throughout (Figure~\ref{fig:atwood}).
\stepn{1} \textbf{Write $\vF=m\va$ for each mass, taking each one's actual direction of motion as
positive.} For $m_1$ (moving \emph{down}): the downward weight $m_1g$ drives it, the upward tension $T$
resists, so $m_1g-T=m_1a$. For $m_2$ (moving \emph{up}): the upward tension $T$ drives it, its weight
$m_2g$ resists, so $T-m_2g=m_2a$.
\stepn{2} \textbf{Add the two equations.} The tension $T$ appears with opposite signs, so it cancels:
$(m_1g-T)+(T-m_2g)=m_1a+m_2a$, giving $(m_1-m_2)g=(m_1+m_2)a$, hence
\[
a=\frac{(m_1-m_2)g}{m_1+m_2}.
\]
\stepn{3} \textbf{Back-substitute for $T$.} Put this $a$ into either equation (say $T=m_2(g+a)$) to get
$T=\dfrac{2m_1 m_2}{m_1+m_2}g$.
\par\medskip\noindent\textbf{Checks.}\ If $m_1=m_2=m$ then $a=0$ and $T=mg$: the system hangs
balanced. If $m_2\to0$ then $a\to g$ and $T\to0$: the lone mass is in free fall. Both limits are
physically correct. Finally, the support holding the pulley carries both string segments, a force $2T$,
and since $(m_1+m_2)^2-4m_1m_2=(m_1-m_2)^2\ge0$,
\[
2T=\frac{4m_1m_2}{m_1+m_2}\,g\le(m_1+m_2)g,
\]
with equality only when $m_1=m_2$. While the masses accelerate, the support carries less than their
total weight.
\end{derivbox}

\subsection{Block on a table, and block on an incline}

Tilting the Atwood machine produces the other two standard machines. The method is unchanged; only the
components of gravity along each body's direction of motion differ.

\begin{figure}[htbp]
\centering
\tikzfig{block_table}
\figcaption{fig:blocktable}{A block on a table joined over a pulley to a hanging mass. The hanging weight
drives the system; friction under the block (if any) opposes it.}
\end{figure}

\begin{derivbox}{Block on a table pulled by a hanging mass}
Block $m_1$ on a horizontal surface (kinetic friction coefficient $\mu_k$) is joined by a light,
inextensible string over a light, frictionless pulley to a hanging mass $m_2$ (Figure~\ref{fig:blocktable}). Suppose the system is
sliding, with $m_2$ descending. Friction is treated fully in \chref{ch:friction}; all that is needed
here is its result that a body sliding on a surface feels a friction force $\mu_kN$ opposing the sliding,
where $N$ is the normal force and $\mu_k$ a dimensionless number for the pair of surfaces (the
coefficient of kinetic friction), and that a body at rest can be held by friction against any push up
to $\mu_sN$, with $\mu_s$ the coefficient of static friction. Vertically, the block's normal force
balances its weight, $N=m_1g$, so kinetic friction is $\mu_k m_1g$.
For $m_1$ (horizontal): $T-\mu_k m_1 g=m_1 a$. For $m_2$ (vertical): $m_2 g-T=m_2 a$. Adding:
\[
a=\frac{(m_2-\mu_k m_1)g}{m_1+m_2},\qquad T=\frac{m_1 m_2(1+\mu_k)g}{m_1+m_2}.
\]
On a smooth table ($\mu_k=0$): $a=\dfrac{m_2 g}{m_1+m_2}$. Starting from rest, the hanging mass sets the
system moving only if $m_2g>\mu_s m_1g$, that is $m_2>\mu_s m_1$. Otherwise the
system stays at rest with $T=m_2g$, and the formulas above, which assume sliding, do not apply.
\end{derivbox}

\begin{figure}[htbp]
\centering
\tikzfig{incline_pulley}
\figcaption{fig:inclinepulley}{The incline version: replace the table's flat surface with a slope, and
the driving term becomes $m_2 g-m_1 g\sin\theta$. The table is just this machine at $\theta=0$; a
vertical Atwood is it at $\theta=90^\circ$.}
\end{figure}

\begin{derivbox}{Incline connected to a hanging mass}
Block $m_1$ sits on a frictionless incline (angle $\theta$), joined by a light, inextensible string over a light, frictionless pulley at the top
to a hanging mass $m_2$ (Figure~\ref{fig:inclinepulley}). Suppose $m_2$ is heavy enough to descend, dragging $m_1$ up the slope; both share
the acceleration magnitude $a$.
\stepn{1} \textbf{For $m_1$, work along the incline.} Perpendicular to the slope nothing accelerates (the
normal force balances $m_1g\cos\theta$), so only the along-slope forces matter: the tension
$T$ pulls it up the slope, while the down-slope component of gravity, $m_1g\sin\theta$, holds it back. So
$T-m_1g\sin\theta=m_1a$.
\stepn{2} \textbf{For $m_2$, work vertically.} Its weight $m_2g$ pulls down, the tension $T$ holds up, and
it accelerates down: $m_2g-T=m_2a$.
\stepn{3} \textbf{Add to cancel $T$:}
\[
a=\frac{(m_2-m_1\sin\theta)g}{m_1+m_2},\qquad T=\frac{m_1 m_2(1+\sin\theta)g}{m_1+m_2}.
\]
Because the incline is smooth, the sign of $a$ tells you which way the system really moves: if it comes
out negative, $m_1$ actually slides \emph{down} and $m_2$ rises, and the same formula still holds. The system balances ($a=0$) exactly when $m_2=m_1\sin\theta$. This is the
master machine: set $\theta=90^\circ$ and it is the Atwood machine; set $\theta=0$ and it is the
block-on-a-table.
\end{derivbox}

\section{Springs}

The ideal string transmits force but never stretches, so its tension is fixed only by the motion. Many
problems instead join bodies by something that stretches measurably, and then the force is fixed by how
far it has stretched. A spring is the simplest such body, and Hooke's law describes it well over a
limited range. As with strings, the springs in this chapter are
\emph{light}: their mass is negligible, so the force is the same at both ends of each spring. Springs
matter beyond springs themselves: a stretched wire, a compressed table top and a bent ruler all behave
like stiff springs for small deformations, so Hooke's law is the simplest model of every elastic force.

\begin{keybox}{Hooke's law}
For deformation $x$ from natural length, the spring force is $F=-kx$: proportional to the deformation and
directed \emph{back} towards the natural length (the minus sign; Figure~\ref{fig:hookes}). Here $x$ is the extension ($x>0$) or
compression ($x<0$), \emph{not} the total length; $k>0$ is the stiffness (or force constant), in
\si{N/m}. The law is empirical and holds only for small deformations, up to the limit of
proportionality.
\end{keybox}

\begin{figure}[htbp]
\centering
\tikzfig{hookes}
\figcaption{fig:hookes}{The spring force is always restoring: whether stretched ($x>0$) or compressed
($x<0$), $F=-kx$ points back towards the natural length. A spring at its natural length exerts no force.}
\end{figure}

\subsection{Series and parallel}

A network of springs that carries a load along one line behaves as a single spring with some effective
stiffness $k_e$. The rule for finding it depends entirely on what the springs share, as Figure~\ref{fig:springs}
shows for two springs.

\begin{derivbox}{Combining springs}
\textbf{Series (springs joined end to end).} Pull the bottom of the chain with force $F$. Because each
spring is light, the net force on it must be zero, so the same $F$ acts across \emph{every} spring in the
line, passed along as tension is along a rope. Each spring then stretches by its own amount,
$x_1=F/k_1$ and $x_2=F/k_2$, and the total stretch is the sum:
\[
x=x_1+x_2=\frac{F}{k_1}+\frac{F}{k_2}=F\left(\frac{1}{k_1}+\frac{1}{k_2}\right).
\]
Comparing with $x=F/k_e$ gives $\dfrac{1}{k_e}=\dfrac{1}{k_1}+\dfrac{1}{k_2}$. Since $1/k_e$ exceeds
both $1/k_1$ and $1/k_2$, $k_e$ is smaller than either, and a series combination is \emph{softer} than
its parts.

\textbf{Parallel (springs side by side).} Now the two springs share the same top and the same bottom
(a bar, loaded so that it stays level), so they stretch by the \emph{same} amount $x$. Each pushes back with its own force, and together they support the
load: $F=k_1x+k_2x=(k_1+k_2)x$. Comparing with $F=k_e x$ gives $k_e=k_1+k_2$, \emph{stiffer} than
either. The test is always: \emph{what is shared?} Shared force $\Rightarrow$ series; shared extension
$\Rightarrow$ parallel.
\end{derivbox}

\begin{figure}[htbp]
\centering
\tikzfig{springs_sp}
\figcaption{fig:springs}{Series springs share the force (extensions add, combination softer); parallel
springs share the extension (forces add, combination stiffer).}
\end{figure}

\begin{trapbox}{springs combine opposite to resistors}
Series springs are \emph{softer} and parallel springs \emph{stiffer}, the reverse of series/parallel
resistors. The reason is that stiffness $k$ behaves like a \emph{conductance} (the reciprocal of
resistance), not a resistance. Reduce
a network by finding a cluster that is plainly in series or plainly in parallel, collapsing it to one
spring, and repeating.
\end{trapbox}

One more property of springs matters whenever something changes suddenly, and it is where springs
differ most from strings.

\begin{keybox}{A spring force cannot jump}
A spring's force is fixed by its deformation, $|F|=k|x|$, and the deformation changes only as the
attached bodies move, which takes time. So at the instant a string is cut or a support removed, every
spring force keeps its previous value. An ideal string or a rigid surface has no such memory: its
tension or normal force can change instantly.
\end{keybox}

\begin{exbox}{Cutting the string above a spring}\label{ex:cutspring}
Block A (mass $m$) hangs from the ceiling by a string, and block B (also $m$) hangs below A from a light
spring. Find the acceleration of each just after the string is cut.

\stepn{1} \textbf{Before the cut.} The spring holds B, so its force is $mg$; the string holds both, $T=2mg$.

\stepn{2} \textbf{Just after.} The tension vanishes at once, but the spring force is still $mg$. On B,
the spring's $mg$ up balances its weight: $a_B=0$. On A, its weight and the spring's pull both act
downwards: $ma_A=2mg$, so $a_A=2g$ downwards.

\stepn{3} \textbf{Check.} The net external force on the pair is $2mg$, and indeed $ma_A+ma_B=2mg$.
\ans{$a_A=2g$; $a_B=0$}
\end{exbox}

\section{Pulley and Spring Combinations}

When strings run over several pulleys, the bodies no longer share one acceleration, and the relation
between their accelerations has to be derived. One procedure does it every time.

\begin{methodbox}{Constraint relations from string length}
\textbf{1.} Measure the position $x_i$ of each moving body or pulley from a fixed reference (the ceiling
or a fixed pulley), along its own line of motion. \textbf{2.} Write the length of each string as the sum
of its straight segments in terms of the $x_i$; the bits wrapped round pulleys add only a constant.
\textbf{3.} Set that length constant and differentiate twice, giving one linear relation among the
accelerations for each string. A minus sign in the result means the bodies move in opposite senses
along the chosen axes.
\end{methodbox}

\subsection{The movable pulley}

A pulley that moves with the load changes both the force needed and the distance moved, and it is the first
system in which the string constraint is not obvious at a glance.

\begin{figure}[htbp]
\centering
\tikzfig{movable_pulley}
\figcaption{fig:movable}{A load $m_2$ on a movable pulley. One string runs from the ceiling, under the movable pulley and over a fixed one to $m_1$. Both segments holding the movable pulley carry the same tension $T$, so the upward pull on the load is $2T$; in exchange the load moves half as far, and accelerates half as fast, as the free end ($a_1=2a_2$).}
\end{figure}

\begin{derivbox}{The movable pulley and its mechanical advantage}
A load $m_2$ hangs from a movable pulley; one ideal string runs from a ceiling anchor, under the movable
pulley, over a fixed pulley, and down to $m_1$ (Figure~\ref{fig:movable}).
\stepn{1} \textbf{Find the constraint from the string length.} The string has a fixed total length. Two
of its segments run down to the movable pulley and back up; if that pulley descends by a distance $d$,
\emph{each} of those two segments must lengthen by $d$, so together they take up $2d$ of string. That
string can only come from the free end, so the free end (carrying $m_1$) must rise by $2d$; equally, if
$m_1$ descends by $2d$, the load rises by $d$. Differentiating twice, the magnitudes of the accelerations
obey the same ratio. (In the language of the method above: with $x_p$ the pulley's depth below the
ceiling and $x_1$ that of $m_1$ below the fixed pulley, the string length is $\ell=2x_p+x_1+\text{constant}$, so
differentiating twice gives $2a_p+a_1=0$ for the downward accelerations $a_p$ of the pulley (and of the
load hanging from it) and $a_1$ of $m_1$: $m_1$ always moves opposite to the pulley, and twice as fast.)
In magnitudes, with $a_2$ the load's acceleration,
\[
a_1=2a_2\qquad(\text{the free end moves, and accelerates, twice as fast as the load}).
\]
\stepn{2} \textbf{Write $\vF=m\va$ for each body.} Suppose $m_1$ descends and the load rises. The string
tension $T$ is the same throughout. The light movable pulley is held up by \emph{two} segments, each
pulling with $T$, so the upward force on the load is $2T$: thus $2T-m_2 g=m_2 a_2$ (upwards positive).
For the descending $m_1$ (downwards positive): $m_1 g-T=m_1 a_1$.
\stepn{3} \textbf{Solve using $a_1=2a_2$.} Substitute $a_1=2a_2$ into the $m_1$ equation, then eliminate
$T$ between the two:
\[
a_2=\frac{(2m_1-m_2)g}{4m_1+m_2},\qquad a_1=2a_2,\qquad T=m_1(g-a_1)=\frac{3m_1m_2\,g}{4m_1+m_2}.
\]
If $2m_1<m_2$, $a_2$ is negative: the load descends and $m_1$ rises.
\stepn{4} \textbf{Mechanical advantage.} To just \emph{hold} the load ($a=0$) needs only $2T=m_2 g$, so
the free-end pull is $T=\tfrac12 m_2 g$, half the weight. The price is distance: you pull twice as far,
so the work done ($\text{force}\times\text{distance}$) is the same as in lifting the load directly. An
ideal machine multiplies force, never work; a real one, with friction, delivers less work than it
receives.
\end{derivbox}

\subsection{Reaching a spring through a pulley}

When a string ends on a spring, the two must agree about the force at the junction, and that agreement is all
that is needed.

\begin{figure}[htbp]
\centering
\tikzfig{spring_pulley}
\figcaption{fig:springpulley}{A spring reached through a fixed pulley. The string carries the spring's force, so $T=kx$ at the junction, and in equilibrium $T=mg$ at the mass: the stretch is $x=mg/k$, exactly as if $m$ hung from the spring directly. A \emph{fixed} pulley only redirects the force and leaves the effective stiffness unchanged.}
\end{figure}

\begin{methodbox}{Coupling a spring to a string}
Wherever a string ends on a spring, the tension \emph{is} the spring force: $T=kx$. For \emph{equilibrium}
set the net force to zero ($T$ fixed by the weights, stretch $x=T/k$). For \emph{motion}, the spring
supplies the restoring force: reduce the system to $m\,d^2x/dt^2=-k_e x$, with $x$ measured from
equilibrium. The motion is then oscillatory with period $2\pi\sqrt{m/k_e}$, a result quoted here because
simple harmonic motion is studied later in the course.
\end{methodbox}

A fixed pulley leaves the stiffness unchanged, as Figure~\ref{fig:springpulley} shows. A movable
pulley does not, and the string constraint says by how much.

\begin{figure}[htbp]
\centering
\tikzfig{spring_movable}
\figcaption{fig:springmovable}{A spring reached through a \emph{movable} pulley. When the load moves down
by $x$, both segments under the pulley lengthen by $x$, and the only place the extra $2x$ of string can
come from is the spring. Twice the stretch gives twice the tension change, and two segments each carry it,
so the load feels four times the stiffness.}
\end{figure}

\begin{derivbox}{A movable pulley multiplies the effective stiffness}
A load hangs from a light movable pulley; the string under it is fixed to the ceiling at one end and, at
the other end, runs up to a spring of stiffness $k$ anchored above (Figure~\ref{fig:springmovable}).
Displace the load downwards by $x$ from equilibrium.
\stepn{1} \textbf{Constraint.} Both string segments lengthen by $x$. The other end of the string is fixed
to the ceiling, so all $2x$ of extra string must come from the spring, which stretches by an extra $2x$.
\stepn{2} \textbf{Force.} The extra tension is therefore $\Delta T=k(2x)$, and the two segments that hold
the pulley each carry it, so the extra upward force on the load is $2\Delta T=4kx$.
\stepn{3} \textbf{Result.} The load feels an effective stiffness $k_e=4k$: the spring acts four times as
stiff. In general, if $n$ string segments hold the load, the spring stretches $n$ times as far as the load
moves and the restoring force is $n$ times the tension, so $k_e=n^2k$. Always read $n$ off the
string-length constraint.
\end{derivbox}

\section{Summary}\label{sec:summary8}
An ideal string has one tension throughout and links the accelerations of the bodies it joins through a
length constraint. With that constraint written down, each body gets its own free-body diagram and its own
equation, and the unknown tension is eliminated between them. Springs obey Hooke's law up to the limit of
proportionality, and a network of springs reduces to one effective stiffness by asking what the springs share.

\begin{itemize}
\item Atwood machine: $a=\dfrac{(m_1-m_2)g}{m_1+m_2}$ and $T=\dfrac{2m_1m_2g}{m_1+m_2}$.
\item The force on a pulley's support is the vector sum of the tensions on either side, not their scalar sum
      unless the strings are parallel.
\item Hooke's law $F=-kx$ uses the deformation $x$, never the total length.
\item Series springs share the same force and are softer: $1/k_e=1/k_1+1/k_2$. Parallel springs share the same
      extension and are stiffer: $k_e=k_1+k_2$. This is the opposite of the rule for resistors.
\item Derive pulley constraints from the string length. A load hung from a movable pulley moves half as far
      as the free end of the string, and a spring at the free end of such a string acts four times as
      stiff on the load ($k_e=n^2k$ when the spring stretches $n$ times as far as the load moves).
\end{itemize}

\section{Problem Types}\label{sec:types8}
Connected-body questions take five recurring forms. Most are solved by the same three moves: write the
string constraint, draw one free-body diagram per body, and add or subtract the resulting equations so that
the tension cancels. When a question asks for the tension as well, substitute the acceleration back into
whichever single-body equation is simplest.

\archq{Two masses over a pulley, find $a$ and $T$}
{``Masses $m_1$ and $m_2$ hang over a light frictionless pulley (or one lies on a table); find the
acceleration and the tension.''}
{Write Newton's law for each body taking its own direction of motion as positive, then \emph{add} the
two equations to eliminate $T$. For a hanging pair, $a=\dfrac{(m_1-m_2)g}{m_1+m_2}$ and
$T=\dfrac{2m_1m_2g}{m_1+m_2}$.}
{Setting $T$ equal to one of the weights. While the masses accelerate, $T$ lies strictly between $m_2g$
and $m_1g$. Also, $T$ is the same on both sides only for a light, frictionless pulley.}
{Masses of \SI{6}{kg} and \SI{4}{kg} hang from a light frictionless pulley ($g=\SI{10}{m/s^2}$). The
tension in the string is:}
{\SI{20}{N}}{\SI{40}{N}}{\SI{48}{N}}{\SI{100}{N}\archans{C}{$T=\dfrac{2m_1m_2g}{m_1+m_2}
=\dfrac{2(6)(4)(10)}{10}=\SI{48}{N}$. Note it lies between $m_2g=\SI{40}{N}$ and $m_1g=\SI{60}{N}$, as
it must; options (B) and (D) are the traps of using one weight or the total weight, and (A) is the
net driving force $(m_1-m_2)g$.}}

\archq{Block on a table connected to a hanging mass}
{``A block on a horizontal table is joined over a pulley to a hanging mass; find $a$ and $T$.''}
{Write Newton's second law for each and add to eliminate $T$:
$a=\dfrac{(m_2-\mu_k m_1)g}{m_1+m_2}$ once the system slides, reducing to $\dfrac{m_2g}{m_1+m_2}$ on a smooth table. Then get $T$
from the \emph{hanging} mass: $T=m_2(g-a)$.}
{Setting $T=m_2g$. That holds only if the system is \emph{not} accelerating; while the hanging mass
accelerates downwards, $T$ is \emph{less} than its weight.}
{A \SI{4}{kg} block on a smooth table is connected to a hanging \SI{6}{kg} mass ($g=\SI{10}{m/s^2}$).
The acceleration of the system is:}
{\SI{4}{m/s^2}}{\SI{6}{m/s^2}}{\SI{10}{m/s^2}}{\SI{15}{m/s^2}\archans{B}{$a=\dfrac{m_2g}{m_1+m_2}
=\dfrac{6(10)}{10}=\SI{6}{m/s^2}$, with $T=6(10-6)=\SI{24}{N}$. Option (A) puts the block's mass in the
numerator, (C) treats the hanging mass as falling freely, and (D) divides by the block's mass alone.}}

\archq{Force on the pulley and its support}
{``What is the force on the pulley, or the reading of a balance holding it?''}
{The two string segments each pull with $T$, so the force on the pulley is the \emph{vector} sum. For two
vertical segments this is $2T$, and for a right-angle bend it is $T\sqrt2$. Compute $T$ first.}
{Answering $(m_1+m_2)g$. That would hold only if nothing accelerated; with unequal masses, $2T$ is
strictly \emph{less} than the total weight.}
{Masses of \SI{5}{kg} and \SI{3}{kg} hang from a light pulley ($g=\SI{10}{m/s^2}$). The force on the
pulley (from the string) is:}
{\SI{75}{N}}{\SI{80}{N}}{\SI{37.5}{N}}{\SI{20}{N}\archans{A}{$T=\dfrac{2(5)(3)(10)}{8}=\SI{37.5}{N}$,
so the pulley carries $2T=\SI{75}{N}$, less than the total weight \SI{80}{N}, because the system
accelerates. Option (B) is the trap above, and (C) counts only one segment.}}

\archq{Springs in series and parallel}
{``Two springs are joined; find the effective force constant, or the extension under a load.''}
{\textbf{Series} (same force, extensions add): $\dfrac{1}{k_e}=\dfrac{1}{k_1}+\dfrac{1}{k_2}$, so
$k_e$ is \emph{smaller} than either. \textbf{Parallel} (same extension, forces add):
$k_e=k_1+k_2$, \emph{larger} than either. Cutting a spring into $n$ equal parts multiplies each
part's $k$ by $n$ (each part carries the same force but stretches only $1/n$ as much).}
{Swapping the two rules; they are the reverse of the resistor rules, which is exactly why they are
mixed up.}
{Springs of \SI{200}{N/m} and \SI{300}{N/m} are joined in \emph{series}. The effective force constant
is:}
{\SI{500}{N/m}}{\SI{250}{N/m}}{\SI{120}{N/m}}{\SI{100}{N/m}\archans{C}{
$\dfrac1{k_e}=\dfrac1{200}+\dfrac1{300}=\dfrac{5}{600}$, so $k_e=\SI{120}{N/m}$, less than both, as series
always is. Option (A) is the parallel rule, the trap above.}}

\archq{Masses on an incline joined over a pulley}
{``One mass lies on an incline, another hangs; find the acceleration.''}
{Resolve the incline mass's weight along the slope ($m_1g\sin\theta$). On a rough incline, first decide
whether anything moves: the net drive $m_2g-m_1g\sin\theta$ must exceed, in magnitude, the maximum static
friction $\mu_s m_1g\cos\theta$. If it does not, the system stays at rest. If it does, kinetic friction
opposes the actual motion; for $m_2$ descending,
$a=\dfrac{m_2g-m_1g\sin\theta-\mu_k m_1g\cos\theta}{m_1+m_2}$, and for $m_1$ sliding down, reverse the
signs of the first two terms. On a smooth incline drop the friction term, and the sign of $a$ gives the
direction.}
{Using $m_1g$ instead of $m_1g\sin\theta$ for the incline mass. Only the along-slope component competes.}
{A \SI{2}{kg} block on a smooth $30^\circ$ incline is joined over a pulley to a hanging \SI{3}{kg} mass
($g=\SI{10}{m/s^2}$). The acceleration is:}
{\SI{2}{m/s^2}}{\SI{4}{m/s^2}}{\SI{6}{m/s^2}}{\SI{10}{m/s^2}\archans{B}{
$a=\dfrac{3(10)-2(10)(0.5)}{5}=\dfrac{30-10}{5}=\SI{4}{m/s^2}$. Option (A) uses $m_1g$ in place of
$m_1g\sin\theta$ (the trap above), and (C) ignores the incline mass's weight altogether.}}

\section{Challenge Problems}\label{sec:challenge8}
The three problems below extend \chsref{ch:newton}{ch:connected}. They are harder than examination questions, chiefly
because each requires the constraint or the free-body diagram to be set up without the cue of a familiar
diagram. Attempt each fully before reading its solution.

\begin{challengebox}{From a chain of blocks to a heavy rope (builds on Newton's laws)}
Three identical \SI{2}{kg} blocks sit in a row on a frictionless table, joined by light strings. A
horizontal force $F=\SI{12}{N}$ pulls the front block.
(a) Find the acceleration and the tension in each string.
(b) Now replace the blocks and strings by a uniform rope of mass $M$ and length $\ell$, lying straight on
the table and pulled at one end by $F$. Find the tension at a distance $x$ from the pulled end.
\challengeans{(a) Treat all three blocks as one system: $a=\dfrac{F}{3m}=\dfrac{12}{6}=\SI{2}{m/s^2}$.
Each string accelerates everything \emph{behind} it. The rear string pulls only the last block:
$T_{\text{rear}}=ma=\SI{4}{N}$. The front string, between the first and second blocks, pulls the last
two: $T_{\text{front}}=2ma=\SI{8}{N}$. Check on the front block: $F-T_{\text{front}}=12-8=\SI{4}{N}=ma$.
(b) The whole rope has $a=F/M$. Isolate the part beyond the point $x$: its mass is $M(\ell-x)/\ell$, and the only
horizontal force on it is the tension $T(x)$ at the cut. Hence
\[
T(x)=\frac{M(\ell-x)}{\ell}\cdot\frac{F}{M}=F\left(1-\frac{x}{\ell}\right),
\]
falling linearly from $F$ at the pulled end to zero at the free end. The staircase of tensions in (a),
$12$, $8$, $4$, $0$ newtons, is the discrete version of this line. A rope with mass does \emph{not} carry
one tension throughout, which is why the single-tension rule of \chref{ch:newton} requires a light
string. \ans{\SI{2}{m/s^2}; \SI{8}{N}, \SI{4}{N}; $F(1-x/\ell)$}}
\end{challengebox}

\begin{challengebox}{The double Atwood machine (builds on pulley constraints)}
A light string over a fixed pulley carries a \SI{3}{kg} mass on its left end. Its right end supports a
light movable pulley, over which a second light string carries masses of \SI{1}{kg} and \SI{3}{kg}. Both
pulleys are frictionless; take $g=\SI{10}{m/s^2}$. The right-hand side holds \SI{4}{kg} in total and the
left only \SI{3}{kg}. Does the left-hand mass rise? Find the acceleration of every mass and both
tensions.
\challengeans{Let the movable pulley have upward acceleration $A$, so the left mass has acceleration $A$
\emph{downwards}. Let the \SI{1}{kg} mass accelerate upwards at $b$ \emph{relative to the pulley}; the
second string has fixed length, so the \SI{3}{kg} mass accelerates at $b$ downwards relative to it. The
light movable pulley needs zero net force, so the main tension is $T_1=2T_2$. With upwards positive:
\[
2T_2-30=-3A,\qquad T_2-10=1(A+b),\qquad T_2-30=3(A-b).
\]
Dividing the last equation by 3 and adding it to the second eliminates $b$:
$\tfrac43T_2-20=2A$. The first gives $A=10-\tfrac23T_2$, so $\tfrac43T_2-20=20-\tfrac43T_2$, hence
$T_2=\SI{15}{N}$, $T_1=\SI{30}{N}$ and $A=0$. Then $b=T_2-10=\SI{5}{m/s^2}$.
The left-hand mass stays at rest, although it is outweighed: the \SI{1}{kg} mass rises and the \SI{3}{kg}
mass falls, each at \SI{5}{m/s^2}. The movable side pulls on the main string with $2T_2=\SI{30}{N}$, not
with its weight of \SI{40}{N}, because its own masses accelerate (the result of the Atwood derivation
above: a support carries less than the total weight while the masses accelerate). In effect the pair
behaves as a single mass $4m_2m_3/(m_2+m_3)=\SI{3}{kg}$ (with $m_2=\SI{1}{kg}$ and $m_3=\SI{3}{kg}$),
exactly balancing the left side.
\ans{left mass at rest; $\pm\SI{5}{m/s^2}$; $T_1=\SI{30}{N}$, $T_2=\SI{15}{N}$}}
\end{challengebox}

\begin{challengebox}{Two springs and a movable pulley (builds on combining springs)}
Two springs, $k_1=\SI{100}{N/m}$ and $k_2=\SI{300}{N/m}$, hang from the ceiling. A light string joins
their lower ends and passes under a light, frictionless movable pulley, from which a \SI{3}{kg} mass hangs;
both string segments are vertical. When the mass is attached and the system has settled, how far has the
mass descended? What is the effective stiffness of the arrangement, as felt by the mass?
($g=\SI{9.8}{m/s^2}$.)
\challengeans{\emph{Force first.} The pulley is held by two segments of one string, so $2T=mg$ and
$T=\SI{14.7}{N}$; both springs carry this same force. Their extensions are $x_1=T/k_1=\SI{0.147}{m}$ and
$x_2=T/k_2=\SI{0.049}{m}$. \emph{Then the constraint.} The two segments together have a fixed length, so
when their upper ends descend by $x_1$ and $x_2$ the pulley descends by the average,
\[
y=\frac{x_1+x_2}{2}=\SI{0.098}{m}\approx\SI{9.8}{cm}.
\]
The effective stiffness is $k_e=mg/y=29.4/0.098=\SI{300}{N/m}$. In symbols, $y=\tfrac{mg}{4}\left(\tfrac1{k_1}
+\tfrac1{k_2}\right)$, so $k_e=\dfrac{4k_1k_2}{k_1+k_2}$: the springs share one tension and so act in
series ($k_s=\SI{75}{N/m}$), and the movable pulley multiplies that by $n^2=4$. Both ideas of the chapter
appear in one answer, and the order matters: identify what is shared and what the string constraint is
before combining anything. \ans{\SI{9.8}{cm}; \SI{300}{N/m}}}
\end{challengebox}
